\documentclass[a4paper,10pt]{article}

\usepackage[utf8]{inputenc}
\usepackage[T1]{fontenc}
\usepackage{geometry}
\usepackage{booktabs}
\usepackage{array}
\usepackage{tabularx}
\usepackage{longtable}
\usepackage{listings}
\usepackage{xcolor}
\usepackage[hidelinks]{hyperref}
\usepackage{amsmath}

\geometry{
    a4paper,
    top=1.8cm,
    bottom=1.8cm,
    left=1.8cm,
    right=1.8cm
}

\setlength{\parindent}{0pt}
\setlength{\parskip}{5pt}
\renewcommand{\arraystretch}{1.18}

\newcommand{\hex}[1]{\texttt{0x#1}}
\newcommand{\code}[1]{\texttt{#1}}

\title{\textbf{DEOS Steering Node \& Service ICD}\\
\large Sürüm 1.1 -- Direksiyon Düğümü (\hex{11}) ve Servisi (\hex{02})}
\author{\textbf{Emir Arkalı} (\href{mailto:emirarkali07@gmail.com}{emirarkali07@gmail.com})\\DEOS -- Dokuz Eylül Otonom Sistemler}
\date{İlk Yayın: 28 Eylül 2026 \quad | \quad Son Revizyon: 13 Ekim 2026}

\begin{document}
\maketitle

\begin{center}
\small
\begin{tabularx}{\textwidth}{ll lX}
\toprule
\textbf{Doküman Tipi:} & Interface Control Document (ICD) & \textbf{Doküman Kodu:} & \code{DOC-ICD-NODE-STR-009} \\
\textbf{Doküman Sürümü:} & Sürüm 1.1 & \textbf{Protokol Sürümü:} & v1.1 (\hex{11}) -- Medya Bağımsız \\
\textbf{Yazar:} & Emir Arkalı (\texttt{emirarkali07@gmail.com}) & \textbf{Son Değiştiren:} & Emir Arkalı \\
\textbf{İlk Yayın Tarihi:} & 28 Eylül 2026 & \textbf{Son Revizyon Tarihi:} & 13 Ekim 2026 \\
\textbf{Durum:} & Yürürlükte / Onaylandı & \textbf{Taşıma / Kapsam:} & Steering Düğümü (\hex{11}) \& STEERING Servisi (\hex{02}) \\
\bottomrule
\end{tabularx}
\end{center}

\subsection*{Doküman Revizyon Geçmişi}
\begin{center}
\small
\begin{tabularx}{\textwidth}{c c l X}
\toprule
\textbf{Revizyon} & \textbf{Tarih} & \textbf{Yazar / Değiştiren} & \textbf{Açıklama} \\
\midrule
v0.3 & 28 Eylül 2026 & Emir Arkalı & İlk taslak ve monolitik ICD v0.3 tanımı. \\
v1.0 & 28 Eylül 2026 & Emir Arkalı & Direksiyon alt sistemi bağımsız doküman revizyonu. \\
v1.1 & 13 Ekim 2026 & Emir Arkalı & Sürüm 1.1 standardizasyonu, düğüm ve servis ayrımı netleştirildi, kurumsal üstveri eklendi. \\
\bottomrule
\end{tabularx}
\end{center}
\vspace{0.8em}

\section{Direksiyon (Steering) Düğüm ve Servis Mimarisi}

DEOS mimarisinde direksiyon alt sistemi, hem donanımsal/mantıksal bir \textbf{Düğüm (Node)} hem de bu düğüm üzerinde çalışan birincil \textbf{Servis (Service)} olarak tanımlanır:

\begin{itemize}
    \item \textbf{Mantıksal Düğüm (Node ID):} \code{DEOS\_NODE\_STEERING} (\hex{11})
    \item \textbf{Birincil Fonksiyonel Servis (Service ID):} \code{DEOS\_SERVICE\_STEERING} (\hex{02})
    \item \textbf{Desteklenen Ortak Servisler:} \code{DEOS\_SERVICE\_SYSTEM} (\hex{00}), \code{DEOS\_SERVICE\_DIAGNOSTIC} (\hex{02})
\end{itemize}

\subsection{Gerçek Zamanlı Kontrol Komutları (Real-Time RX)}
\begin{center}
\begin{tabularx}{\textwidth}{>{\raggedright\arraybackslash}p{3.8cm} >{\raggedright\arraybackslash}p{3cm} >{\raggedright\arraybackslash}p{2.2cm} >{\raggedright\arraybackslash}p{2.2cm} X}
\toprule
\textbf{Değişken Adı} & \textbf{Sınıf \& Servis} & \textbf{Komut ID} & \textbf{Ölçek} & \textbf{Fonksiyonel Görev} \\
\midrule
\code{target\_steering\_angle} & COMMAND (0x1) STEERING (0x02) & 0x01 & 0.01 deg & Hedef tekerlek açısı (Negatif: Sol, Pozitif: Sağ). \\
\code{target\_angular\_velocity} & COMMAND (0x1) STEERING (0x02) & 0x01 & 0.01 deg/s & Direksiyonun hedef açıya döneceği açısal hız tavanı. \\
\code{steering\_control\_mode} & COMMAND (0x1) STEERING (0x02) & 0x04 & Enum & Direksiyon motorunun kontrol modunu belirler. \\
\code{emergency\_stop\_cmd} & SAFETY (0x0) SYSTEM (0x00) & 0x05 & Boolean & Direksiyon motorunu mekanik serbest bırakma (free-wheel). \\
\bottomrule
\end{tabularx}
\end{center}

\subsection{Direksiyon Durum Telemetrisi (STATUS)}

\code{DEOS\_CMD\_STEERING\_GET\_STATUS (0x02)} komutuna karşılık dönen veri yapısı:

\begin{center}
\begin{tabular}{c c c l}
\toprule
\textbf{Offset} & \textbf{Boyut} & \textbf{Tip} & \textbf{Alan} \\
\midrule
0 & 2 byte & int16\_t  & Current Angle (0.01 deg) \\
2 & 2 byte & uint16\_t & Motor Current (0.1 A) \\
4 & 2 byte & int16\_t  & Motor Temperature (0.1 $^\circ$C) \\
6 & 2 byte & int16\_t  & Chassis Yaw Rate (0.01 deg/s) \\
8 & 1 byte & uint8\_t  & Is Tracking Target (0: No, 1: Yes) \\
\bottomrule
\end{tabular}
\end{center}

\subsection{Çapraz Sistem Dinamik Telemetrisi}
\begin{itemize}
    \item \textbf{current\_vehicle\_speed} (0x10): Hız arttıkça azami direksiyon açısını kısarak yüksek hızda devrilmeyi önler.
    \item \textbf{chassis\_yaw\_rate} (0x03): Araç savrulma anında ESP mantığı ile karşı direksiyon (counter-steering) desteği sağlar.
\end{itemize}

\subsection{Direksiyon Konfigürasyon Parametreleri}
\begin{center}
\begin{tabularx}{\textwidth}{>{\raggedright\arraybackslash}p{3.8cm} >{\raggedright\arraybackslash}p{3cm} >{\raggedright\arraybackslash}p{2.2cm} >{\raggedright\arraybackslash}p{2.2cm} X}
\toprule
\textbf{Parametre} & \textbf{ID} & \textbf{Tip} & \textbf{Ölçek} & \textbf{Açıklama} \\
\midrule
PID\_KP / KI / KD & \hex{0001-03} & uint16\_t & 0.01 & Açı kontrolü PID kazançları. \\
PID\_I\_MAX & \hex{0004} & uint16\_t & 0.01 \% & İntegratör doyum patlamasını (windup) önler. \\
FEEDFORWARD\_K & \hex{0005} & uint16\_t & 0.01 & Hedef açısal hız ileri besleme kazancı. \\
DEADBAND & \hex{0006} & uint16\_t & 0.01 deg & Düz gidişte titreşimleri kesen ölü bölge. \\
ZERO\_OFFSET & \hex{0010} & int16\_t & 0.01 deg & Mekanik düz gidiş merkez açı sapması (Kalibrasyon). \\
MIN/MAX\_ANGLE & \hex{0011-12} & int16\_t & 0.01 deg & Yazılımsal sol ve sağ azami dönüş sınırları. \\
CURRENT\_LIMIT & \hex{0013} & uint16\_t & 0.1 A & Direksiyon motoru elektriksel akım tavanı. \\
MAX\_ANGULAR\_VEL & \hex{0014} & uint16\_t & 0.01 deg/s & İzin verilen azami açısal dönüş hızı tavanı. \\
MAX\_ANGULAR\_ACC & \hex{0015} & uint16\_t & 0.01 deg/s² & Mekanik dişliyi koruyan açısal ivme limiti. \\
SOFT\_STOP\_MARGIN & \hex{0016} & uint16\_t & 0.01 deg & Mekanik dayamaya sert çarpmayı önleyen pay. \\
MAX\_TRACK\_ERR & \hex{0017} & uint16\_t & 0.01 deg & Takip hatası toleransı; aşılırsa hata verilir. \\
TRACK\_TIMEOUT\_MS & \hex{0018} & uint16\_t & 1 ms & İzleme hatasının sürdürülebileceği azami süre. \\
\bottomrule
\end{tabularx}
\end{center}

\section{Kalibrasyon Prosedürleri}
Direksiyon sistemi \code{DEOS\_SERVICE\_CALIBRATION (0x05)} servisini kullanır.
\begin{itemize}
    \item \textbf{Direksiyon:} Lazerli terazi ile tam düz konum alınarak enkoder açısı \code{ZERO\_OFFSET} olarak EEPROM'a yazılır. Motor düşük torkla sola ve sağa dayamalara sürülerek nominal akım eşiğinde durduğu noktalar yazılımsal sınır açıları (\code{MIN\_ANGLE}, \code{MAX\_ANGLE}) atanır.
\end{itemize}

\end{document}
